\chapter{Conical structure and convergence of supports}\label{sec:varifold-conical-structure}
This section supplies the homogeneity argument in tangent-cone construction and separates weak measure convergence from convergence of supports. The source for the cone argument is \cite{Minter}, Lemmas 1.1--1.2, pp.~11--12. Integral compactness remains the external input in Theorem \ref{thm:allard-compactness}.

\begin{lemma}[Stationarity is closed under varifold convergence]\label{lem:stationarity-closed}
If stationary $m$-varifolds $V_j$ converge weakly to $V$ in $U$, then $V$ is stationary in $U$.
\end{lemma}
\begin{proof}
For $X\in C_c^1(U;\mathbb R^N)$ the function $(x,S)\mapsto\operatorname{div}_S X(x)$ is continuous and compactly supported on the Grassmann bundle. Thus $\delta V(X)=\lim_j\delta V_j(X)=0$. This assertion alone does not guarantee that an arbitrary limit is integral; integrality uses Theorem \ref{thm:allard-compactness} under its hypotheses.
\end{proof}

\begin{lemma}[The weight determines a rectifiable varifold]\label{lem:rectifiable-lift-unique}
Two rectifiable $m$-varifolds with the same weight are equal as measures on the Grassmann bundle.
\end{lemma}
\begin{proof}
Write their weights as $\theta_i\mathcal H^m\llcorner M_i$ with positive densities. Equality of the weights implies that $M_1$ and $M_2$ differ by an $\mathcal H^m$-null set, after discarding null subsets. Their approximate tangent planes therefore agree almost everywhere by the rectifiable tangent characterization, Theorem \ref{thm:rectifiable-tangents}. Their Radon--Nikodym densities also agree. Substitute into Definition \ref{def:rectifiable-varifold} for every compactly supported continuous test function. The positivity of the densities and rectifiability are essential to this argument.
\end{proof}

\begin{proposition}[Radial tangency gives homogeneous weight]\label{prop:radial-tangency-homogeneous-weight}
Let $C$ be a stationary $m$-varifold on all of $\mathbb R^N$, with $x\in S$ for $C$-almost every $(x,S)$. Its weight $\nu$ satisfies
\[
 s^{-m}\int\phi(x/s)\,d\nu(x)=\int\phi(x)\,d\nu(x)\qquad(s>0)
\]
for every $\phi\in C_c^1(\mathbb R^N)$. In particular $\nu(sA)=s^m\nu(A)$ for every Borel set $A$.
\end{proposition}
\begin{proof}
Test stationarity with $X(x)=\phi(x)x$. Radial tangency gives
$\int(m\phi+x\cdot D\phi)\,d\nu=0$.
The derivative of $F(s)=s^{-m}\int\phi(x/s)\,d\nu$ is
\[
 F'(s)=-s^{-m-1}\int\bigl(m\phi(x/s)+(x/s)\cdot D\phi(x/s)\bigr)\,d\nu=0,
\]
using the same identity with the test function $x\mapsto\phi(x/s)$. Differentiation is justified by a common compact support and bounded derivatives when $s$ ranges in a compact subinterval of $(0,\infty)$. Thus $F(s)=F(1)$. Equality on smooth tests identifies the two Radon measures and extends to Borel sets.
\end{proof}

\begin{theorem}[Constant mass ratio characterizes rectifiable stationary cones]\label{thm:constant-ratio-cone}
For a stationary rectifiable $m$-varifold $C$ on $\mathbb R^N$, the following are equivalent: $r^{-m}\|C\|(B_r(0))$ is constant for all $r>0$; $x\in S$ for $C$-almost every $(x,S)$; and $(\eta_{0,s})_\#C=C$ for every $s>0$. Compare \cite{Minter}, Lemma 1.2, pp.~11--12.
\end{theorem}
\begin{proof}
The annular identity, Proposition \ref{prop:stationary-annular-identity}, applied on countably many annuli exhausting $\mathbb R^N\setminus\{0\}$, proves radial tangency from the constant ratio. The origin has zero weight by Corollary \ref{cor:stationary-density-defect}. Radial tangency gives the homogeneous weight identity by Proposition \ref{prop:radial-tangency-homogeneous-weight}. The dilation has tangential Jacobian $s^{-m}$, so the weights of $C$ and $(\eta_{0,s})_\#C$ coincide. Both varifolds are rectifiable; Lemma \ref{lem:rectifiable-lift-unique} identifies their entire lifts. Finally dilation invariance gives $\|C\|(B_r)=r^m\|C\|(B_1)$ directly. Homogeneity of the weight alone would not justify identifying the Grassmannian distributions for a general nonrectifiable varifold.
\end{proof}

\begin{proposition}[Uniform mass lower bound at integral support points]\label{prop:stationary-support-mass-bound}
If $V$ is stationary and integral in $U$, $a\in\operatorname{spt}\|V\|$, and $\overline{B_r(a)}\subset U$, then
\[
 \|V\|(B_r(a))\geq\omega_m r^m.
\]
In particular a nonzero stationary integral varifold on all of $\mathbb R^N$ has infinite total mass.
\end{proposition}
\begin{proof}
Corollary \ref{cor:integral-support-density} gives density at least one. Monotonicity bounds each normalized mass ratio below by its limit at zero. For an entire nonzero varifold choose one support point and let $r\to\infty$.
\end{proof}

\begin{proposition}[Local convergence of stationary integral supports]\label{prop:stationary-support-convergence}
Suppose stationary integral $m$-varifolds $V_j$ converge to a stationary integral varifold $V$ in $U$. For every compact $K\subset U$ and $\varepsilon>0$, for all sufficiently large $j$,
\[
 \operatorname{spt}\|V_j\|\cap K\subset
 \{x:\operatorname{dist}(x,\operatorname{spt}\|V\|)<\varepsilon\},
\]
\[
 \operatorname{spt}\|V\|\cap K\subset
 \{x:\operatorname{dist}(x,\operatorname{spt}\|V_j\|)<\varepsilon\}.
\]
Distances are to the full supports in $U$, not their intersections with $K$; the distance to an empty support is infinite.
\end{proposition}
\begin{proof}
If the first inclusion fails, choose support points $a_j\in K$ staying at least $\varepsilon$ from the limiting support, and pass to $a_j\to a\in K$. Choose $r>0$ small enough that $\overline{B_{2r}(a)}\subset U$ misses that support. A nonnegative compactly supported cutoff equal to one on $B_r(a)$ then has integrals against $\|V_j\|$ tending to zero. But $B_{r/2}(a_j)\subset B_r(a)$ for large $j$, so Proposition \ref{prop:stationary-support-mass-bound} gives a fixed positive lower bound, a contradiction.

If the second inclusion fails, choose $a_j\in\operatorname{spt}\|V\|\cap K$ tending to $a$ while $B_\varepsilon(a_j)$ misses $\operatorname{spt}\|V_j\|$. The relative support is closed, so $a$ belongs to it. A fixed sufficiently small ball about $a$ lies in these empty balls for large $j$. All compactly supported test integrals in that ball vanish in the limit, contradicting that $a$ is a support point. Only the first direction needs the uniform integral mass lower bound.
\end{proof}
